% Hình 1.1: Minh họa vòng đời lỗ hổng phân quyền có trạng thái (Stateful Broken Access Control - A1)
% Thiết kế đơn sắc chuẩn học thuật: Khoảng cách giữa Victim và Gateway được mở rộng lên 6.2cm,
% đảm bảo tuyệt đối không có nhãn chữ nào bị tràn hay đè qua đường gióng thời gian.
\begin{tikzpicture}[
  >=Stealth,
  scale=0.88, transform shape,
  actor/.style={rectangle, draw=black, fill=gray!12, thick, minimum width=2.4cm, minimum height=0.9cm, align=center, rounded corners=2pt, font=\bfseries\small},
  server/.style={rectangle, draw=black, fill=gray!12, thick, minimum width=2.6cm, minimum height=0.9cm, align=center, rounded corners=2pt, font=\bfseries\small},
  db/.style={cylinder, draw=black, fill=gray!12, shape border rotate=90, aspect=0.25, thick, minimum width=2.4cm, minimum height=1.0cm, align=center, font=\bfseries\small},
  timecol/.style={font=\bfseries\footnotesize, anchor=east}
]

  % 1. Bốn thực thể tham gia (Participants) xếp ngang ở đỉnh
  \node[actor]  (victim) at (0, 0)      {Nạn nhân};
  \node[server] (api)    at (6.2, 0)    {Máy chủ REST API};
  \node[actor]  (admin)  at (10.4, 0)   {Quản trị viên};
  \node[db]     (db)     at (14.6, 0)   {Cơ sở dữ liệu\\người dùng và vai trò};

  % 2. Các đường gióng thời gian (Lifelines)
  \draw[dashed, black!40] (victim) -- (0, -11.0) coordinate (v_end);
  \draw[dashed, black!40] (api)    -- (6.2, -11.0) coordinate (a_end);
  \draw[dashed, black!40] (admin)  -- (10.4, -11.0) coordinate (ad_end);
  \draw[dashed, black!40] (db)     -- (14.6, -11.0) coordinate (d_end);

  % ==================== BƯỚC 1: Login & Thao tác ban đầu ====================
  \node[timecol] at (-0.4, -1.2) {Thời điểm $t_0, t_1$};
  
  \draw[->, thick] (0, -1.2) -- 
    node[midway, above=2pt, font=\scriptsize, fill=white, inner sep=1.5pt] {1. Gửi lời gọi ghi nhạy cảm (token vai trò quản trị)} (6.2, -1.2);
  
  \draw[->, thick] (6.2, -1.9) -- 
    node[midway, above=2pt, font=\scriptsize, fill=white, inner sep=1.5pt] {Xác thực quyền} (14.6, -1.9);
  
  \draw[<-, thick] (6.2, -2.6) -- 
    node[midway, below=2pt, font=\scriptsize, fill=white, inner sep=1.5pt] {Hợp lệ (vai trò quản trị)} (14.6, -2.6);
  
  \draw[<-, thick] (0, -3.3) -- 
    node[midway, above=2pt, font=\scriptsize, fill=white, inner sep=1.5pt] {Phản hồi: \textbf{200 OK}} (6.2, -3.3);

  % ==================== BƯỚC 2: Hạ quyền (Transition) ====================
  \node[timecol] at (-0.4, -4.5) {Thời điểm $t_2$};
  
  \draw[->, thick] (10.4, -4.5) -- 
    node[midway, above=2pt, font=\scriptsize, align=center, fill=white, inner sep=1.5pt] {2. Hạ nạn nhân xuống vai trò thấp hơn} (6.2, -4.5);
  
  \draw[->, thick] (6.2, -5.2) -- 
    node[midway, above=2pt, font=\scriptsize, fill=white, inner sep=1.5pt] {Cập nhật vai trò mới} (14.6, -5.2);
  
  \draw[<-, thick] (6.2, -5.9) -- 
    node[midway, below=2pt, font=\scriptsize, fill=white, inner sep=1.5pt] {Xác nhận đã cập nhật (200 OK)} (14.6, -5.9);

  % ==================== BƯỚC 3: Tái sử dụng token cũ ====================
  \node[timecol] at (-0.4, -7.0) {Thời điểm $t_3$};
  
  \draw[->, thick] (0, -7.0) -- 
    node[midway, above=2pt, font=\scriptsize, fill=white, inner sep=1.5pt] {3. Gửi lại lời gọi ghi bằng \textbf{token cũ}} (6.2, -7.0);

  % Hộp điểm nghẽn bảo mật (nằm trọn vẹn trong khoảng cách 6.2cm, cách lề mỗi bên 0.6cm)
  \node[rectangle, draw=black, very thick, dashed, fill=white, rounded corners=3pt, 
        align=center, font=\scriptsize, inner sep=5pt, text width=5.0cm] at (3.1, -8.3) {
    \textbf{Nguyên nhân lỗ hổng:}\\
    máy chủ chỉ đọc vai trò trong token hoặc phiên đã lưu,\\
    không tra lại vai trò hiện hành trong cơ sở dữ liệu.
  };

  % Kết quả thực tế (Lỗ hổng)
  \draw[<-, very thick] (0, -9.6) -- 
    node[midway, above=2pt, font=\bfseries\scriptsize, fill=white, inner sep=1.5pt] {Lỗ hổng: vẫn trả về \textbf{200 OK}} (6.2, -9.6);

  % Kỳ vọng an toàn (Dashed)
  \draw[<-, thick, dashed] (0, -10.4) -- 
    node[midway, above=2pt, font=\itshape\scriptsize, fill=white, inner sep=1.5pt] {Kỳ vọng: từ chối với \textbf{403 Forbidden}} (6.2, -10.4);

\end{tikzpicture}
